% 第6章第8节：本章小结。

\section{本章小结}
\label{sec:graph-spectral-summary}

本章沿“顶点值—边差—加权汇总—谱坐标—滤波或划分”一路推进。关联矩阵的转置把顶点函数变成边差，按边权缩放后再汇总回顶点，就得到组合拉普拉斯$\bm L=\bm B\bm W_E\bm B^{\top}=\bm D-\bm W$，其二次型是沿边变化能量。连接总量不均时，度加权内积与坐标缩放给出$\bm L_{\mathrm{sym}}$，邻域平均给出$\bm L_{\mathrm{rw}}$，表~\ref{tab:graph-spec-laplacian-comparison}汇总了三者的差别。谱定理把这些算子的特征向量排成从平缓到剧烈的模态：零特征值的重数等于正权支撑分量数，$\bm L_{\mathrm{sym}}$的谱落在$[0,2]$中。在此基础上要分清三组区别：沿边变化能量按特征值加权，幅值平方不加权；图傅里叶变换只换坐标，滤波才改变函数；谱松弛给出下界和候选方向，离散分组必须经阈值扫描恢复，并用原目标重新计算。

\begin{table}[!htb]
  \centering\small
  \begin{tabularx}{\linewidth}{@{}lXXX@{}}
    \toprule
    & $\bm L$ & $\bm L_{\mathrm{sym}}$ & $\bm L_{\mathrm{rw}}$\\
    \midrule
    正度部分的定义 & $\bm D-\bm W$ & $\bm D^{-1/2}\bm L\bm D^{-1/2}$ & $\bm I-\bm D^{-1}\bm W$\\
    比较的量 & 相邻顶点的原值之差 & 除以$\sqrt{d_i}$后的值之差 & 本点值与邻居平均之差\\
    对称性 & 欧氏内积下对称 & 欧氏内积下对称 & 度加权内积下自伴\\
    每个分量的零模 & 指示向量$\bm1_C$ & $\bm D^{1/2}\bm1_C$；孤点为$\bm e_v$ & 指示向量$\bm1_C$\\
    谱的范围 & $[0,2d_{\max}]$，随边权尺度变化 & $[0,2]$ & 与$\bm L_{\mathrm{sym}}$相同\\
    支撑孤点 & 行列为零 & 约定$\bm D^{-1/2}$该处为$0$，行列为零 & 约定$(\bm P_{\mathrm{rw}})_{v,v}=1$，该行为零\\
    \bottomrule
  \end{tabularx}
  \caption{同一非负权图上三种拉普拉斯的对照。“分量”指正权支撑图的连通分量。$\bm L_{\mathrm{sym}}$与$\bm L_{\mathrm{rw}}$在正度顶点上相似，特征值相同，特征向量相差因子$\bm D^{1/2}$；三者的特征向量不能直接互换。}
  \label{tab:graph-spec-laplacian-comparison}
\end{table}

回看整个部分：第~\ref{chap:graph-objects}章建立了图、路径、连通与矩阵表示；第~\ref{chap:graph-traversal}章用遍历提取连通结构、拓扑次序和结构标记；第~\ref{chap:graph-optimization}章在路径、生成树、流和匹配上求最优；第~\ref{chap:graph-set-functions}章把评价对象扩展到子集；第~\ref{chap:graph-dynamic-programming}章利用局部依赖结构计算有限配置。前五章的对象都是离散的，本章转向挂在顶点上的数值，增加的是图算子与特征坐标这两样工具。它们并不取代前面的算法：判断连通仍可以直接遍历，最小割仍有精确的流算法，谱方法只在自己的条件下给出能量、坐标、滤波和松弛下界。随机游走的概率解释、连续扩散的演化和完整的傅里叶计算留到后续部分接续；本章所用的谱理论始终停留在有限维线性代数的范围内。
